\subsection{A Rule of Life for Advent}\label{sec:rule}

A rule of life is a short written plan of what one will do, each day and each week, to keep a
season. It turns good intentions into a few definite acts and makes them easy to review. The rules
below are the book's proposal. Only the cells marked \textit{Law} state an obligation, and each of
them names the law it comes from; every other cell is counsel.

The three columns are three measures of the same rule. \textit{Beginning} suits someone who has
not kept Advent before, or whose duties leave little room. \textit{Steady} suits someone who
already prays daily and wants the season to shape the day. \textit{Demanding} suits someone with
more freedom and some experience of the spiritual life, and it should be taken up only with a
confessor's or director's advice. Most readers should choose the first or second column, keep it
whole, and leave the third alone this year.

\medskip
{\small
\begin{longtable}{@{}>{\raggedright\arraybackslash}p{0.14\linewidth}>{\raggedright\arraybackslash}p{0.25\linewidth}>{\raggedright\arraybackslash}p{0.25\linewidth}>{\raggedright\arraybackslash}p{0.25\linewidth}@{}}
\toprule
Practice & Beginning & Steady & Demanding\\
\midrule
\endfirsthead
\toprule
Practice & Beginning & Steady & Demanding\\
\midrule
\endhead
Prayer &
\textit{Counsel:} the Angelus morning, noon and evening; the \latin{Alma Redemptoris Mater} at night. &
\textit{Counsel:} as at the beginning, with Vespers or Compline of the Little Office in its Advent
form. &
\textit{Counsel:} as steady, with the hymn \latin{Creator alme siderum} at evening and the O
antiphon of each day from 17 to 23 December.\\
\addlinespace
Scripture &
\textit{Counsel:} the Sunday Epistle and Gospel, read on the Saturday evening. &
\textit{Counsel:} the Isaiah lesson of each day (Appendix~\ref{sec:daily}). &
\textit{Counsel:} the day's Isaiah read in the four steps of reading, meditation, prayer and
contemplation (Section~\ref{sec:prayer}).\\
\addlinespace
Mass &
\textit{Law} (Latin Church, Code of 1983, cc.~1246--1247, current 2026): Mass on Sundays and on
holy days of obligation, among them 8 December unless the conference of bishops has suppressed it
or moved it to a Sunday (Section~\ref{sec:second-sunday}). &
\textit{Counsel:} one weekday Mass each week, and Mass on one of the Ember days. &
\textit{Counsel:} daily Mass where it can be had.\\
\addlinespace
Fasting &
\textit{Law} (Latin Church, Code of 1983, cc.~1250--1251, current 2026): every Friday is a
penitential day, with abstinence unless it is a solemnity. \textit{Law} (United States, bishops'
norms of 1966, pars.~22--25, kept in force 1983): Friday remains penitential; abstinence from meat
has first place among its forms, but outside Lent is not binding under sin. &
\textit{Counsel:} as at the beginning, with the three Ember days kept as days of fast. (The Code
of 1917 required this from 1918 to 1983; neither the Code of 1983 nor the United States norms of
1966 requires it.) &
\textit{Counsel:} as steady, with the Wednesdays and Fridays of Advent kept with fast and
abstinence, as the English custom noted in 1870 kept them, and only with a director's advice.\\
\addlinespace
Alms &
\textit{Counsel:} one fixed gift to the poor, given before Christmas. &
\textit{Counsel:} the money the fast saves, given away each week, as Leo asks. &
\textit{Counsel:} as steady, with one work of mercy done in person each week.\\
\addlinespace
Silence and examen &
\textit{Counsel:} the examen at night in Ignatius's five points. &
\textit{Counsel:} as at the beginning, with one hour each Sunday free of news and entertainment. &
\textit{Counsel:} as steady, with a fixed time of silence each day and St Alphonsus's meditation
for the day (Appendix~\ref{sec:daily}).\\
\addlinespace
Confession &
\textit{Counsel:} confession once before Christmas. &
\textit{Counsel:} confession in Ember week. &
\textit{Counsel:} confession early in Advent and again in Ember week or before the Vigil.\\
\bottomrule
\end{longtable}
}

\subsubsection{Keeping a rule}

Write the rule down on the evening before the First Sunday, on one side of a card, and keep it
where you pray. Read it once a week, on the Saturday evening, and ask three questions of the week
just past: what was kept, what was dropped, and why. The examen of each night
(Section~\ref{sec:prayer}) does the same work day by day.

A rule fails in two ways, and both are common. The first is too much: a plan drawn up in the
fervor of the First Sunday that collapses by the second week and leaves the season worse than if
nothing had been attempted. Choose less than you think you can do. The second is drift: the rule is
kept for a few days and then quietly forgotten. When that happens, take it up again the next
morning without fuss, and go on from there.

Do not raise the rule in the middle of the season. If the beginning column proves easy, keep it
faithfully to Christmas and take up the steady column next year. If something must change, let it
be because of a duty to others, an illness or a confessor's word, and change it downward.

\subsubsection{Households}

Much of the rule can be kept together. The Angelus at the evening meal, the O antiphon of the day
read aloud in the last week, the Sunday Gospel read on Saturday night, a shared gift to the poor:
these suit a household as well as a single person, and children can take part in all of them.
Fasting is different. Under the Code of 1983 minors are not bound by the law of fast, nor those who
have not completed their fourteenth year by the law of abstinence, and the same Code asks pastors and parents to see that they are formed in a true
sense of penance all the same (can.~1252). A
household fast that suits the adults should be explained to the children and adapted for them, not
imposed.

\subsubsection{Sundays and feasts}

The Sundays of Advent are of the first class, and the rule may be relaxed on them. The Code of
1917 suspended the law of fast on Sundays (can.~1252 \S\,4), and a private fast is best set aside
on them too. The feast of the Immaculate Conception on 8 December is kept
as a feast. In 2026 the feasts of St Andrew (30 November) and St Thomas (21 December) fall on
weekdays; a household may mark them, and the St Andrew prayer begins on the first of them
(Section~\ref{sec:prayer}).
